\chapter{Blacher, Boris} % (fold)
\label{cha:blacher}

\href{https://en.wikipedia.org/wiki/Boris_Blacher}{Composer page} on Wikipedia.

\bigskip



\noindent{Visions fugitives, Op. 22 (1915–17, arranged 1935)} \\
\noindent{Habemeajaja, chamber opera (1929, premiered 1987)} \\
\noindent{Jazz-Koloraturen, Op. 1 for coloratura soprano, alto saxophone and bassoon (1929)} \\
\noindent{String Quartet No. 1, Op. 11 (1930)} \\
\noindent{String Trio - Three Studies on Jewish folksongs (1931)} \\
\noindent{Fünf Sinnsprüche Omars des Zeltmachers, Op. 3 for mezzo-soprano and piano (1931)} \\
\noindent{Konzert-Ouvertüre (1931)} \\
\noindent{Two Toccatas for piano (1931)} \\
\noindent{Kleine Marschmusik, Op. 2 for orchestra (1932)} \\
\noindent{Orchester-Capriccio über ein Volkslied, Op. 4 for orchestra (1933)} \\
\noindent{Alla Marcia for orchestra (1934)} \\
\noindent{Divertimento for string orchestra (1935)} \\
\noindent{Etüde for string quartet (1935)} \\
\noindent{Fest im Süden, Op. 6, Danse-drama in one act by Ellen Petz (1935)} \\
\noindent{Fest im Süden - Suite for orchestra (1935)} \\
\noindent{Divertimento for wind orchestra, Op. 7 (1936)} \\
\noindent{Concertante Musik, Op. 10 for orchestra (1937)} \textit{Rundfunk-Sinfonieorchester Berlin, Johannes Kalitzke}\marginnote{high} \\
\noindent{Symphony, Op. 12 (1938)} \\
\noindent{Dance Scenes La Vie, ballet in one act (1938)} \\
\noindent{Rondo for orchestra (1938)} \\
\noindent{Harlekindae, Op. 13, ballet in one act with prologue and epilogue by Jens Keith (1939)} \\
\noindent{Fürstin Tarakanowa (Princess Tarakanova), Op. 19, opera in three acts (1940)} \\
\noindent{Fürstin Tarakanowa, Op. 19a - Suite for orchestra (1940)} Overture: \textit{Rundfunk-Sinfonieorchester Berlin, Johannes Kalitzke}\marginnote{high} \\
\noindent{String Quartet No. 2, Op. 16 (1940)} \\
\noindent{Flute Sonata, Op. 15 (1940)} \\
\noindent{Concerto for String Orchestra, Op. 20 (1940)} \\
\noindent{Hamlet; Symphonic Poem, Op. 17 for orchestra (1940)} \textit{Rundfunk-Sinfonieorchester Berlin, Johannes Kalitzke}\marginnote{high} \\
\noindent{Two Sonatinas, Op. 14 for piano (1940)} \\
\noindent{Cello Sonata (1940–41)} \\
\noindent{Violin Sonata, Op. 18 (1941)} \\
\noindent{Sonatina for Piano Four Hands (1942)} \\
\noindent{Der Großinquisitor (The Grand Inquisitor), Op. 21 - Oratorio after Dostoyevsky by Leo Borchard (1942)} \\
\noindent{Three Pieces for piano, Op. 23 (1943)} \\
\noindent{Romeo und Julia, Op. 22, chamber opera, premiered Salzburg Festival (1950)} \\
\noindent{Drei Psalmen for baritone and piano (Psalms 142, 141 \& 121), arranged for chamber ensemble by Beyer (1943)} \\
\noindent{Piano Sonata No. 1, Op. 14/1 (1943)} \\
\noindent{Piano Sonata No. 2, Op. 14/2 (1943)} \\
\noindent{String Quartet No. 3, Op. 32 (1944)} \\
\noindent{Vier Chöre on texts by François Villon (1944)} \\
\noindent{Partita for strings and 6 percussion, Op. 24a (1945)} \\
\noindent{Concerto for Jazz Orchestra (1945)} \\
\noindent{Chiarina, Op. 33 - ballet in one act by Paul Strecker (1946)} \\
\noindent{Die Flut, Op. 24b, radio opera (1946)} \\
\noindent{Divertimento for trumpet, trombone, and piano, Op. 31 (1946)} \\
\noindent{Die Flut (The Tide), chamber opera in one act (text by Heinz von Cramer) (1946)} \\
\noindent{Piano Concerto No. 1, Op. 28 (1947)} \\
\noindent{Orchestral Variations on a Theme by Paganini, Op. 26 (1947)} \textit{Wiener Philharmoniker, Sir Georg Solti}\marginnote{high} \\
\noindent{Die Nachtschwalbe (The Night Swallow), Op. 27, Zeitoper in one act (1947)} \\
\noindent{Four Songs, Op. 25 for soprano and piano. Text by Friedrich Wolf (1947)} \\
\noindent{Violin Concerto, Op. 29 (1948)} \\
\noindent{Hamlet, Op. 35 - ballet in a prologue and three scenes after Shakespeare by Tatjana Gsovsky (1949)} \\
\noindent{Hamlet - Suite for orchestra (1949)} \\
\noindent{Preußisches Märchen (A Prussian Fairytale), ballet-opera in six scenes (1949/52)} \\
\noindent{Ornamente, Sieben Studien über variable Metren, Op. 37 for piano (1950)} \\
\noindent{Lysistrata, Op. 34 ballet in three scenes after Aristophanes (1950)} \\
\noindent{Lysistrata - Suite, Op. 34a from ballet for orchestra (1950)} \textit{Rundfunk-Sinfonieorchester Berlin, Johannes Kalitzke}\marginnote{high} \\
\noindent{Concerto for Clarinet, Bassoon, Horn, Trumpet and Harp, Op. 36 (1950)} \\
\noindent{Piano Sonata No. 3, Op. 39 (1951)} \\
\noindent{Divertimento for four woodwinds, Op. 38 (1951)} \\
\noindent{Epitaph: In memory of Franz Kafka, (String Quartet No. 4), Op. 41 (1951)} \\
\noindent{Sonata for solo violin, Op. 40 (1951)} \\
\noindent{Dialog for flute, violin, piano, and string orchestra (1951)} \\
\noindent{Nebel, for voice and piano (1951)} \\
\noindent{Piano Concerto No. 2 (in variable metres), Op. 42 (1952)} \\
\noindent{Orchester-Ornament, Op. 44 (1953)} \\
\noindent{Studie im Pianissimo, Op. 45 for orchestra (1953)} \\
\noindent{Abstrakte Oper Nr. 1, Op. 43, experimental opera in one act (1953/57)} \\
\noindent{Two Inventions, Op. 46 for orchestra (1954)} \\
\noindent{Viola Concerto, Op. 48 (1954)} \\
\noindent{Francesca da Rimini, Op. 47 - fragment from Dante's Divina Commedia, for soprano and solo violin (1954)} \\
\noindent{Der Mohr von Venedig, Op. 50 - ballet in 6 scenes and an epilogue after Shakespeare by Erika Hanka (1955)} \\
\noindent{Träume vom Tod und vom Leben, Op. 49 - Cantata after a poem by Hans Arp for tenor, choir and orchestra (1955)} \\
\noindent{Orchester-Fantasie, Op. 51 (1956)} \\
\noindent{Hommage à Mozart - Metamorphoses on a group of Mozart themes, for orchestra (1956)} \\
\noindent{Music for Cleveland, Op. 53 for orchestra (1957)} \\
\noindent{Two Poems, Op. 55 for jazz quartet (1957)} \\
\noindent{Thirteen Ways of Looking at a Blackbird, Op. 54 for soprano and string quartet (or piano) (1957/64)} \\
\noindent{Aprèslude, Op. 57 for voice and piano (1958)} \\
\noindent{Songs of the Sea Pirate O'Rourke, Op. 56 for solo voices and orchestra (1958)} \\
\noindent{Requiem, Op. 58 (1958)} \\
\noindent{Musica giocosa, Op. 59 for orchestra (1959)} \\
\noindent{Rosamunde Floris, Op. 60 - opera (libretto by Gerhart von Westerman, based on the play by Georg Kaiser) (1960)} \\
\noindent{Jüdische Chronik (A Jewish Chronicle), for chorus and orchestra (1961)} \\
\noindent{Variations on a theme of Muzio Clementi, Op. 61 for piano and orchestra (
